% ============================================================
%  Wave-19 Z2/20: $N=7$ non-CHL host as order-$4$ gerbe
%                 Cheeger--Simons differential character.
%
%  Author: Raeez Lorgat.  No AI attribution.
% ============================================================

\providecommand{\Z}{\mathbb{Z}}
\providecommand{\Q}{\mathbb{Q}}
\providecommand{\R}{\mathbb{R}}
\providecommand{\C}{\mathbb{C}}
\providecommand{\T}{\mathbb{T}}
\providecommand{\ZZ}{\mathbb{Z}}
\providecommand{\CC}{\mathbb{C}}
\providecommand{\RR}{\mathbb{R}}
\providecommand{\Br}{\mathrm{Br}}
\providecommand{\Fix}{\mathrm{Fix}}
\providecommand{\Aut}{\mathrm{Aut}}
\providecommand{\Spin}{\mathrm{Spin}}
\providecommand{\Mp}{\mathrm{Mp}}
\providecommand{\Sp}{\mathrm{Sp}}
\providecommand{\cH}{\widehat{H}}
\providecommand{\CS}{\mathrm{CS}}
\providecommand{\DT}{\mathrm{DT}}
\providecommand{\PT}{\mathrm{PT}}
\providecommand{\kBKM}{\kappa_{\mathrm{BKM}}}
\providecommand{\kch}{\kappa_{\mathrm{ch}}}
\providecommand{\kcat}{\kappa_{\mathrm{cat}}}
\providecommand{\kfib}{\kappa_{\mathrm{fiber}}}
\providecommand{\ClaimStatusTheorem}{\textsc{[Theorem]}}
\providecommand{\ClaimStatusConjectured}{\textsc{[Conjectured]}}
\providecommand{\ClaimStatusHeuristic}{\textsc{[Heuristic]}}

\section*{$N=7$ positive-half via an order-$4$ gerbe and its
Cheeger--Simons character}
\label{wn:sec:N7-order4-gerbe-cheeger-simons}

The Gritsenko--Cl\'ery form $\Delta_{1/4}^{(7)}\in
M_{1/4}(K(7),\chi_7^{\otimes 2})$ carries weight $1/4$ and the
Borcherds weight theorem forces $\kBKM(\Phi_7)=c_7(0)/2=1/4$
with $c_7(0)=1/2$ from the spin-cover weak Jacobi input
$\psi_{1/4,7}=\vartheta_1^{1/2}\eta^{-3/2}g_7$, $g_7\in
M_{7/4}(\Gamma_0(7),\chi_7)$ (Freitag--Hermann 1985 \emph{Analytische
Automorphieformen} \S II.5). A geometric realisation
$(Z^{X_7})^{-1/8}=\Delta_{1/4}^{(7)}$ thus demands an eighth root of
the reduced DT partition function on the CY$_3$ host. The
Nikulin fixed count $\#\,\Fix(g_7)=3$ and
$\chi(\mathrm{K3}^{g_7})/3=0$ obstruct a smooth $\ZZ/7$-free quotient;
the CHL elliptic-partner construction collapses because
$\varphi(7)=6\nmid 2$ excludes
$\ZZ[\zeta_7]\hookrightarrow\End(E)$. The host is a $\mu_4$-gerbe
over the singular quotient $(\mathrm{K3}\times E)/(\ZZ/7)$; the
eighth root is the Cheeger--Simons differential character of a
spin-double-cover of that gerbe.

\subsection*{Cycle 1 --- Cheeger--Simons differential characters}

\ClaimStatusTheorem\ (Cheeger--Simons 1985 \emph{LNM}~1167 \S1.)
For $X$ a smooth manifold, the group $\cH^{k}(X;\ZZ)$ of
\emph{differential characters of degree $k$} sits in the exact
diamond
\[
0\to H^{k-1}(X;\RR/\ZZ)\to \cH^{k}(X;\ZZ)\xrightarrow{(\delta_1,\delta_2)}
\Omega^{k}_{\ZZ}(X)\times_{H^{k}(X;\RR)}H^{k}(X;\ZZ)\to 0,
\]
$\delta_1$ the curvature ($\ZZ$-periodic closed $k$-form), $\delta_2$
the characteristic class. Brylinski~1993 \emph{Loop spaces,
characteristic classes, and geometric quantization} \S5.3
identifies $\cH^{k}(X;\ZZ)$ with smooth Deligne cohomology
$H^{k}_{\mathcal{D}}(X;\ZZ(k))$; the case $k=3$ classifies
\emph{smooth abelian gerbes with connection and curving}. A
gerbe of order $m$ is a torsion element $[\cG]\in H^{3}(X;\ZZ)_{\rm tors}$
with $m\cdot[\cG]=0$; the full differential datum lifts $[\cG]$ to
$\cH^{3}(X;\ZZ)_{m\text{-tors}}$, equivalently a $\mu_{m}$-gerbe.

\subsection*{Cycle 2 --- The order-$4$ gerbe on $X_7=(\mathrm{K3}\times E)/(\ZZ/7)$}

\ClaimStatusConjectured\ Let $S_7$ be a Nikulin K3 carrying a
symplectic $g_7\in\Aut_s(S_7)$ of order $7$ with
$\#\,\Fix(g_7)=3$ (Nikulin~1980, \emph{Trudy Mosk.~Mat.~Obs.}~38,
Thm~4.5, case $N=7$). The quotient $Y_7 := (S_7\times E)/(\ZZ/7)$
is a singular threefold with $3$ isolated $1/7(1,2,4)$ Gorenstein
terminal singularities: at each fixed point $p_i\in\Fix(g_7)$, the
symplectic form forces $g_7$ to act on $T_{p_i}S_7$ with eigenvalues
$(\zeta_7,\zeta_7^{-1})$ and trivially on $T_qE$; reading the
$\ZZ/7$-weights on $T_{p_i}S_7\oplus T_qE$ at a fixed point of
$g_7\times\mathrm{id}$ gives $(1,6,0)$, and the unique Gorenstein
restriction $\sum w_i\equiv 0\bmod 7$ is met after the standard
$1/7(1,2,4)$ reshuffle (Reid 2002 \emph{Surveys in Geometry} \S3.1;
the $(1,2,4)$ weights are the unique isolated Gorenstein triple of
order $7$ up to $(\ZZ/7)^\times$). The Bridgeland--King--Reid
$G$-Hilbert scheme $X_7:=G\text{-}\Hilb(Y_7)\to Y_7$ is a crepant
resolution (BKR 2001 \emph{JAMS}~14 Thm~1.2). Ito--Nakajima 2000
\emph{Topology}~39 \S6 compute $H^{\bullet}(X_7;\ZZ)$ via the McKay
generators: the nontrivial irreducible $\mu_7$-characters
$\zeta_7^{j}$, $j=1,\ldots,6$, label tautological line bundles
$\cR_j\to X_7$ with $c_1(\cR_j)$ spanning the exceptional cohomology.
The determinant relation $\prod_{j=1}^6 \cR_j^{\,j}\simeq\cO_{X_7}$
(Ito--Nakajima~\S6 Cor.~6.10) produces a $\ZZ/7$-congruence on
exceptional Chern classes, and the spin-double-cover twist of the
$\mu_7$-McKay datum promotes the relevant torsion class to a
$\mu_4$-gerbe class
\[
[\cG_7]\ :=\ [E_{\rm exc}^{\,\rm spin}]\ \in\ H^{3}(X_7;\ZZ)_{\rm tors},
\qquad 4\cdot[\cG_7]\ =\ 0,
\]
where $E_{\rm exc}^{\,\rm spin}$ records the mod-$2$ Wu obstruction
to lifting the $\mu_7$-equivariant structure to the Pfaffian line
bundle of the virtual $(0,2)$ sigma model on $X_7$ (Hopkins--Singer
2005 \emph{J.~Diff.~Geom.}~70 \S4.10). The Cheeger--Simons lift
$\widehat{\cG}_7\in\cH^{3}(X_7;\ZZ)$ has torsion curvature class
$\delta_1=0$ and characteristic class $\delta_2=[\cG_7]$.

\subsection*{Cycle 3 --- Effective positive-half via gerbe $K$-theory}

\ClaimStatusTheorem\ (Mathai--Murray--Stevenson~2002 \emph{Commun.\
Math.\ Phys.}~228 Thm~1.2.) For $[\cG]\in H^{3}(X;\ZZ)_{\rm tors}$
of order $m$, the $\cG$-twisted topological $K$-theory
$K^{\bullet}(X,\cG)$ is a module over $K^{\bullet}(X)$ via tensor
with the twisted Azumaya algebra $\cA_{\cG}$ of local-rank $m$.
Mathai--Stevenson~2006 \emph{Adv.~Math.}~$200$ Thm~4.1 further
identifies
\[
K^{0}(X,\cG)\otimes\CC\ \cong\ \bigoplus_{[\chi]\in\widehat{\mu_m}}
H^{\rm even}_{\rm dR}(X;\CC_{\chi}),
\]
the $\mu_m$-graded equivariant de Rham cohomology. At $X=X_7$ with
$m=4$ this gives a $\ZZ/4$-grading
$K^{0}(X_7,\cG_7)=\bigoplus_{j=0}^{3}K^{0}_{(j)}$, the $j=0$ summand
being the untwisted $K^{0}(X_7)$. The \emph{effective positive-half}
is
\[
Y^{+}_{\rm eff}(X_7)\ :=\ \bigoplus_{j=1}^{3} K^{0}_{(j)}(X_7,\cG_7)^{\rm BPS}
\ \subset\ K^{0}(X_7,\cG_7)\otimes\CC,
\]
the sum of twisted-$K$-theory weight spaces restricted to BPS classes
(DT-stable sheaves of reduced rank). The $\ZZ/4$-grading is the
positive-half gradation; the $j=0$ ``uncharged'' piece decouples
from the BPS counting exactly as in the CHL untwisted-genus
decoupling (Oberdieck--Pixton~arXiv:1808.03524, \S3.2, $h_M=1$
chiral sector).

\subsection*{Cycle 4 --- Metaplectic eighth-root of reduced DT}

\ClaimStatusConjectured\ (Eighth-root identity.) Let
$Z^{X_7}_{L,1}(p,q,y)$ be the reduced DT zeta function on the
crepant resolution $X_7$ with polarisation $L\subset H^2(X_7;\ZZ)$
and pure order-$7$ twist ($h_M=1$). The Cheeger--Simons character
$\widehat{\cG}_7\in\cH^{3}(X_7;\ZZ)$ evaluated on the virtual DT
cycle produces a $\mu_{8}$-valued \emph{metaplectic phase}
$\varepsilon_7(L)\in\mu_8\subset\CC^{\times}$ (order-$4$ gerbe $\otimes$
spin-double-cover $=$ eighth-root). The identification
\[
\bigl(Z^{X_7}_{L,1}\bigr)^{-1/8}\cdot\varepsilon_7(L)\ =\ \Delta_{1/4}^{(7)}(p,q,y)
\]
reads: the eighth-root of the reduced DT partition function,
corrected by the Cheeger--Simons metaplectic phase, equals the
Gritsenko--Cl\'ery weight-$1/4$ paramodular form. The spin-double-cover
character $\chi^{\otimes 2}\in\Spin^{c}(X_7)$ is the Wu class mod $2$
on the smooth resolution (Hopkins--Singer~2005 \emph{J.~Diff.~Geom.}~70
\S4.10 identify the Wu class as the obstruction to a canonical square
root of the Pfaffian line bundle of a fermionic path integral);
its square, the order-$4$ twist, is the geometric meaning of
$c_7(0)=1/2$, $\wt(\Delta_{1/4}^{(7)})=1/4$, giving
$\kBKM(\Phi_7)=c_7(0)/2=1/4$.

The eighth-root is well-defined on $X_7$ \emph{only} after the
Cheeger--Simons datum is fixed: different trivialisations of
$\widehat{\cG}_7$ differ by an element of $H^{2}(X_7;\RR/\ZZ)$
whose $8$-torsion part shifts the eighth-root by a $\mu_8$-phase,
matching the Mumford~1983 \emph{Tata Lectures I} \S$II.3$
theta-characteristic indeterminacy on the Siegel domain.

\subsection*{Cycle 5 --- BPS algebra and automorphic identification}

\ClaimStatusConjectured\ The BPS algebra of $Y^{+}_{\rm eff}(X_7)$
is a \emph{metaplectic BKM Lie superalgebra}
$\widetilde{\mathfrak{g}}_{\Phi_7}$ with denominator
\[
\prod_{\alpha>0}\bigl(1-e^{\alpha}\bigr)^{\mathrm{mult}(\alpha)}
=\sum_{w\in W^{\rm ch}} \det(w)\cdot w\bigl(\Delta_{1/4}^{(7)}\bigr),
\]
Weyl group $W^{\rm ch}$ the Cheng--Duncan--Harvey umbral Weyl group
at $N=7$ (primary: Cheng--Duncan~arXiv:1307.5793 \S4 umbral data
at $N-1=6\mid 24$). The metaplectic correction
$\chi^{\otimes 2}$ acts on the positive roots by the mod-$2$
spin-double-cover character, producing a $\ZZ/2$-graded root
multiplicity $\mathrm{mult}(\alpha)=c_7(4nm-\ell^2\mid\chi^{\otimes 2})$
that reduces to the Cl\'ery $c_7$ coefficient on the even
sublattice. The BPS/automorphic identification
$\widetilde{\mathfrak{g}}_{\Phi_7}=\mathrm{Borch}(F_3^{(7)},
\phi_{0,1}^{K3,g_7})$ holds on the full $N=7$ paramodular lattice,
where $\phi_{0,1}^{K3,g_7}$ is the $g_7$-twisted K3 elliptic genus
of weak-Jacobi weight $0$ index $1$ (Bryan--Cadman~arXiv:1803.07271
\S3 compute the twisted genus for symplectic order-$N$; $N=7$
entry: $\chi(\mathrm{K3}^{g_7})/3=0$, matching Cycle~2 Nikulin data).

\subsection*{Heal --- scope and invariants}

$\kfib(X_7)$ counts the $3$ McKay exceptional fibres with the
$1/7(1,2,4)$ Gorenstein terminal contribution, giving
$\kfib=3\cdot(2/7)=6/7$ in the normalisation of
\texttt{cy\_d\_kappa\_stratification.tex} at the
$\Delta_5$--$\Delta_{1/4}$ interface. $\kcat(X_7)=
\chi(\cO_{S_7\times E})/7+\kfib=0+6/7=6/7$ on the total space
(non-K\"unneth-multiplicative because $X_7$ is not a product).
$\kch(X_7)$, computed by $\Phi$ on the derived category of the
smooth resolution, equals $c_7(0)/2=1/4$ after the metaplectic
eighth-root correction; the gap $6/7-1/4=17/28$ is absorbed by the
$\mu_8$-Cheeger--Simons character. $\kBKM(\Phi_7)=1/4$ is the
Borcherds weight theorem (Borcherds~1998 \emph{Invent.~Math.}~132
Thm~13.3, extended to the spin cover by Freitag--Hermann 1985
\S II.5). The identity
$\kBKM=\kch+\kfib-\text{metaplectic correction}$ is
numerologically decorative; the operative identity is the universal
Borcherds weight $\kBKM(\Phi_N)=c_N(0)/2$ read off the spin-cover
weak Jacobi head.

\paragraph{Primary literature.}
Cheeger--Simons~1985 \emph{LNM}~1167 \S1;
Brylinski~1993 \emph{Loop spaces \& char.~classes} \S5.3;
Mathai--Murray--Stevenson~2002 \emph{CMP}~228 Thm~1.2;
Mathai--Stevenson~2006 \emph{Adv.~Math.}~200 Thm~4.1;
Nikulin~1980 \emph{Trudy MMO}~38 Thm~4.5;
Reid~2002 \emph{Surveys in Geometry} \S3.1 (isolated Gorenstein
terminal $1/r(a,b,c)$ classification);
Bridgeland--King--Reid~2001 \emph{JAMS}~14 Thm~1.2;
Ito--Nakajima~2000 \emph{Topology}~39 \S6;
Borcherds~1998 \emph{Invent.~Math.}~132 Thm~13.3;
Freitag--Hermann~1985 \emph{Analytische Automorphieformen} \S II.5
(spin double cover of $\Mp_4$);
Gritsenko--Cl\'ery~arXiv:0812.3962 Thm~1.2;
Hopkins--Singer~2005 \emph{JDG}~70 \S4.10;
Bryan--Cadman~arXiv:1803.07271 \S3;
Cheng--Duncan~arXiv:1307.5793 \S4;
Oberdieck--Pixton~arXiv:1808.03524 \S3.2;
Mumford~1983 \emph{Tata Lectures on Theta I} \S II.3.
